\documentclass[11pt]{article}
\usepackage{tikz}
\usepackage{graphicx}
\usepackage[margin=1.0in]{geometry}
\usepackage{etoolbox}
\usepackage{amsmath}

\newcommand{\lessonheader}[2]{%
\begin{center} \Large {MA103: Data-Driven Modeling\\Lesson #1: #2}  \end{center}
\hrule
}

\newcommand{\objectiveitem}[1]{\item #1}
\newcommand{\lessonobjectives}[1]{%
\vspace{-0.1in}
\section*{Lesson Objectives}
\vspace{-0.1in}
\begin{enumerate}
\forcsvlist{\objectiveitem}{#1}
\end{enumerate}
\hrule
}

\newcommand{\subheader}[1]{%
\vspace{-0.1in}
\noindent
\section*{#1}
\vspace{-0.1in}
}

\newcommand{\subsubheader}[1]{%
\vspace{-0.1in}
\noindent
\subsubsection*{#1}
\vspace{-0.1in}
}

\begin{document}
\lessonheader{38}{Objective Functions I}

\lessonobjectives{%
  {Construct an objective function and constraints from context and available data (e.g., costs, capacities, resource limits), with correct units.},
  {Formulate a complete LP algebraically.}%
}

\subheader{Warm-Up: Formulating Constraints}

\begin{enumerate}

\item Record the general formulation for a linear programming model. Which part represents the objective function? Which part represents the constraints?

\vfill

\item Record the general formulation for a linear programming model \textit{using matrix vector notation}.

\vfill

\item You branched Cyber. Congratulations! 

Three years later, you are the mission element lead for a National Mission Team. 
Your team of operators and analysts must both 1) train to conduct Offensive Cyberspace Operations and 2) execute those operations as a part of the Cyber Mission Force.

You have 2 fully qualified operators and 3 fully qualified analysts.
Each operator or analyst can work a maximum of 60 hours per week.

Each hour of operations you schedule requires 3 operator-hours and 2 analyst-hours to complete. 
Each hour of training you schedule requires 1 operator-hour and 3 analyst-hours to complete. 
Your higher headquarters requires you to conduct at least 20 hours of operations each week.

Let $x_1$ be the number of hours of operations you schedule each week and $x_2$ be the number of hours of training you schedule each week.
Formulate the constraints using function notation. 

\vfill

\item Formulate the constraints using matrix vector notation.

\vfill

\end{enumerate}

\newpage

\hrule

\subheader{Exercise: Developing Objective Functions}

\subsubheader{Maximization}

\leavevmode\hspace*{\parindent}Imagine you want to maximize the number of hours your team spends on operations. 

Formulate your objective function in function notation and matrix vector notation. 

\vfill

Now you want to prioritize increasing your operational capacity in the long term by increasing your number of qualified operators. 

Qualifying as a Cyberspace Operator requires both training sessions and operational experience. 
You have examined the rate at which your team members progress towards qualification and found that you gain .04 new operators for each hour of operations you schedule and .015 new operators for each hour of training you schedule. 

Formulate an objective function maximizing your team's progress towards qualification.

\vfill

\subsubheader{Minimization}

\leavevmode\hspace*{\parindent}There has been a lot of complaining in your office about the busy new schedule. 

For every hour of operations you schedule, you hear 32 minutes of complaining. 
For every hour of training you schedule, you hear 14 minutes of complaining. 

Formulate an objective function minimizing the amount of complaining on your team.

\vfill

Some of the complaining that you hear comes from fellow team leads, since your operations have been hogging all of the workstations in the Joint Operations Center. 
Eager to rebuild your relationships with your peers, you want to minimize your use of the shared workstations. 

For every hour of operations that you schedule, you use 3 hours of workstation time, since your operator, analyst, and mission commander each need a workstation. 
For every hour of training that you schedule, you reduce the amount of workstation time used by 45 minutes, since your team members' friends are too busy come by and hang out. 

Formulate an objective function minimizing the amount of workstation time used by your team.

\vfill

\newpage

\subsubheader{$>$2 Independent Variables}

\leavevmode\hspace*{\parindent}You want to spend as close to 2/3 of the unit's time on operations as possible without going over.
Formulate an objective function for this goal. Hint: you may need to add an additional constraint.

\vfill

You realize that you have forgotten to account for time dedicated to unit PT. You must conduct a minimum of 5 hours of unit PT each week. PT hours do not count towards the maximum number of hours an individual Soldier can work. Therefore, your new constraints are: 
\begin{align*}
2x_1 + 3x_2 &\le 180 \\
3x_1 + x_2 &\le 120 \\
x_1 &\ge 20 \\
x_2 &\ge 0 \\
x_3 &\ge 5
\end{align*}
Does this affect your ability to graphically depict the objective function?
Does it affect your ability to find an optimal solution to the linear programming model?

\vfill

\subsubheader{Abstract Goals}

\leavevmode\hspace*{\parindent}All of your analysts have been offered jobs at Microsoft. 
You are concerned about retention rates and want to come up with an index that reflects how weekly schedules impact job satisfaction. 

Use your best judgment to develop an objective function that maximizes job satisfaction. Use operations, training sessions, and morning PT as independent variables.

\vfill

Do operations, training sessions, and morning PT contribute positively or negatively to job satisfaction in your objective function? Is the relationship linear? In real life, might these variables have a non-linear relationship with job satisfaction?

\vfill

\hrule

\subheader{Wrap-Up: Food for Thought}

\noindent In the objective functions that we used today, the cost coefficient values were given to you. In the real world, will that always be true? How might you use predictive modeling to calculate parameters for your objective function?

\end{document}